% ============================================================
\chapter{ROD and Comparative Architecture Study}
\label{ch:rod_extension}
% ============================================================

\section{Motivation for Pretrained Feature Transfer}

The PIDNet-S analysis exposed a recurring pattern: low-contrast corridors
were sometimes removed, broad fields could be opened incorrectly, and small
boundary displacements survived even in otherwise accurate masks. With only
1,002 training images, one possible remedy is to reuse visual structure learned
from a much larger corpus rather than asking every feature extractor to learn
from CaT alone. ROD provides a concrete instance of this strategy. It connects
a trainable off-road segmentation decoder to a pretrained EfficientSAM ViT-S
encoder whose weights remain fixed during task training
~\cite{sun2025rod,xiong2024efficientsam}.

This design occupies a markedly different operating point from PIDNet-S. It
introduces substantially more representation capacity and transfers pretrained
features, but the full transformer still executes for every frame. Fixed
weights reduce optimisation cost; they do not make inference inexpensive.
Accuracy and latency must therefore be measured together.

The comparison proceeds from controlled pairwise evidence to architectural
context. PIDNet-S and ROD first share the same split, seeds, epoch budget,
objective, and external resolution. Seven additional systems are then trained
under the same protocol. This distinguishes a gain over one reference model
from a gain over the available compact alternatives.

\section{Implemented ROD ViT-S Model}

The ROD implementation contains 29,106,050 parameters. Of these, 22.35 million
belong to the EfficientSAM ViT-S encoder and are frozen; approximately 6.75
million belong to the decoder and are trainable. Freezing is enforced in three
ways: encoder parameters have \texttt{requires\_grad=False}, the encoder is
kept in evaluation mode even when the outer model enters training mode, and
its forward pass is executed under \texttt{no\_grad}. Batch-normalisation in
the decoder remains trainable.

The external input is the same $3\times384\times640$ tensor used by the other
models. Inside ROD it is bilinearly resized to $1024\times1024$ and normalised
with ImageNet statistics. A $16\times16$ patch embedding produces a
$64\times64$ grid of 384-dimensional tokens. The encoder has 12 transformer
blocks with six attention heads. The learned position embedding was pretrained
on a $14\times14$ patch grid and is bicubically interpolated to the $64\times64$
grid.

The released implementation is more specific than the paper's high-level
description. Rather than summing every transformer output, the decoder uses
the outputs of blocks 2 through 12. Each selected feature is
rearranged to channel-first form and projected from 384 to 128 channels by a
$1\times1$ convolution followed by a $3\times3$ convolution. The projected
features are fused sequentially through residual convolutional blocks. A
post-fusion block refines the result and expands it to 256 channels.

Two residual upsampling stages create progressively larger feature maps. The
decoder then resizes four 256-channel sources---the two upsampled maps, the
pre-upsample map, and the encoder's final image embedding---to one common
resolution, concatenates them, and reduces the result to 256 channels. A
$1\times1$ prediction layer produces two logits. They are finally resized to
the original $640\times384$ external input size.

\begin{figure}[H]
\centering
\resizebox{\textwidth}{!}{\begin{tikzpicture}[
  node distance=0.5cm,
  block/.style={draw,rounded corners=2pt,minimum height=1.15cm,text width=2.25cm,
    align=center,font=\small,inner sep=4pt},
  arrow/.style={-{Latex[length=2.4mm]},thick,draw=azulUC3M}
]
\node[block,fill=softblue] (rgb) {\textbf{RGB tensor}\\$3\times384\times640$};
\node[block,fill=softorange,right=of rgb] (norm) {\textbf{Encoder input}\\resize $1024\times1024$; ImageNet normalise};
\node[block,fill=softpurple,right=of norm] (vit) {\textbf{Frozen ViT-S}\\$16^2$ patches; 12 blocks; 384 channels};
\node[block,fill=softgreen,right=of vit] (latent) {\textbf{Features}\\blocks 2--12 + 256-channel embedding};
\node[block,fill=softorange,below=1.15cm of latent] (decoder) {\textbf{Trainable decoder}\\project, sequentially fuse, residual upsample};
\node[block,fill=softgreen,left=of decoder] (multi) {\textbf{Multiscale fusion}\\four 256-channel scales};
\node[block,fill=softblue,left=of multi] (head) {\textbf{Binary head}\\two logits; resize to $640\times384$};
\node[block,fill=softpurple,left=of head] (mask) {\textbf{Decision}\\softmax and threshold};
\draw[arrow] (rgb) -- (norm);
\draw[arrow] (norm) -- (vit);
\draw[arrow] (vit) -- (latent);
\draw[arrow] (latent) -- (decoder);
\draw[arrow] (decoder) -- (multi);
\draw[arrow] (multi) -- (head);
\draw[arrow] (head) -- (mask);
\node[draw,dashed,fit=(vit)(latent),inner sep=7pt,label={[font=\footnotesize]above:22.35M fixed parameters}] {};
\node[draw,dashed,fit=(decoder)(multi)(head),inner sep=7pt,label={[font=\footnotesize]below:6.75M trainable parameters}] {};
\end{tikzpicture}
}
\caption[Implemented ROD ViT-S pipeline.]{Implemented ROD ViT-S pipeline. The
pretrained encoder is frozen during CaT training but remains part of every
inference. The decoder uses intermediate blocks 2--12 and multiscale
convolutional fusion before the two-class head.}
\label{fig:rod_architecture}
\end{figure}

Figure~\ref{fig:rod_architecture} distinguishes the frozen feature extractor
from the decoder that is updated on CaT.

Full encoder fine-tuning was an available alternative. It was not adopted
because ROD's defining transfer strategy holds the encoder fixed, the CaT
training set is small, and fine-tuning would greatly increase training memory
and the number of task-specific degrees of freedom. The experiment consequently
does not determine whether end-to-end EfficientSAM adaptation would improve the
result.

\section{Controlled PIDNet--ROD Comparison}

The matched comparison uses seeds 1337, 2027, and 4242 under both blue-only and
blue+green label policies. Within a policy, PIDNet-S and ROD receive:
\begin{itemize}
  \item the same 1,002/266/544 manifest assignment;
  \item the same $640\times384$ external tensors and batch size 4;
  \item 15 epochs and validation-mIoU checkpoint selection;
  \item AdamW at $2.5\times10^{-4}$ with weight decay $10^{-4}$;
  \item epoch-level cosine annealing to $10^{-5}$;
  \item the same deterministic augmentation recipe;
  \item biased inverse-frequency class weights, CE+Dice, and threshold 0.50.
\end{itemize}

Matching these decisions prevents a favourable ROD optimiser, seed, or split
from explaining the comparison. It does not make the two networks identical
experiments in every respect: ROD uses pretrained frozen encoder weights,
model-specific normalisation, and an internal $1024\times1024$ transformer input,
whereas PIDNet-S is randomly initialised and operates directly on the external
resolution. The result therefore compares the complete model pipelines under
common task-level controls.

Population mean and population standard deviation are reported:
\begin{align}
 \bar m &= \frac{1}{3}\sum_{k=1}^{3}m_k,\\
 \sigma_{\mathrm{pop}}
 &=\sqrt{\frac{1}{3}\sum_{k=1}^{3}(m_k-\bar m)^2}.
\end{align}
Here, $m_k$ is the value of the reported metric in run $k$; the index
$k\in\{1,2,3\}$ corresponds to seeds 1337, 2027, and 4242; $\bar m$ is their
arithmetic mean; and $\sigma_{\mathrm{pop}}$ is their population standard
deviation. The divisor is three because the equations describe the complete
declared set of runs rather than estimating variance from a random sample.
With only three selected seeds, the standard deviation describes observed
repeatability; it is not a confidence interval over all possible training
runs.

\section{Reconstructed ROD Training Regime}

The controlled recipe answers a fair architecture-comparison question, but it
does not use the optimisation choices documented for ROD. A second experiment
therefore transfers the published settings as closely as the available
description permits. Table~\ref{tab:rod_recipe_config} contrasts the two
recipes.

\begin{table}[H]
\centering
\caption{Training choices held in the controlled architecture comparison and
their counterparts in the closest documented ROD recipe.}
\label{tab:rod_recipe_config}
\small
\begin{tabularx}{\textwidth}{@{}lXX@{}}
\toprule
\textbf{Setting} & \textbf{Controlled comparison} & \textbf{ROD paper approximation} \\
\midrule
Objective & weighted cross-entropy + Dice & cross-entropy only \\
Class weights & inverse frequency with 2.5/0.75 biases & neutral $[1,1]$ \\
AdamW learning rate & $2.5\times10^{-4}$ & $10^{-3}$ \\
Weight decay & $10^{-4}$ & 0.01 \\
Batch size & 4 & 8 (validation 4) \\
Schedule & epoch-level cosine to $10^{-5}$ & iteration-level polynomial, power 0.9 \\
Augmentation & common tested recipe enabled & disabled because none is documented \\
Epoch budget & 15 & 15 retained because publication does not specify one \\
Threshold & 0.50 & 0.50 \\
\bottomrule
\end{tabularx}
\end{table}

The polynomial scheduler updates after every optimiser step:
\begin{equation}
 \eta_t=\eta_0\left(1-\frac{t}{T}\right)^{0.9},
\end{equation}
where $\eta_t$ is the learning rate after optimiser step $t$, $\eta_0=10^{-3}$
is the initial learning rate, $T$ is the total number of optimiser steps in
15 epochs, and $0.9$ is the fixed polynomial-decay exponent transferred from
the documented recipe. The step index runs from $t=0$ to $t=T$. The run uses
bfloat16 mixed precision, the official EfficientSAM ViT-S checkpoint, and a
frozen encoder. Because the ROD paper does not disclose its epoch budget,
augmentation policy, or all timing details, this experiment is called the
\emph{closest documented recipe}. It is evidence about transfer of the
reported settings to CaT, not an exact reproduction of a published CaT result.

\section{Architecture Set}

After the PIDNet--ROD comparison, seven additional compact systems were added:
FPN/MobileNetV2, FPN/EfficientNet-B0, U-Net/MobileNetV2,
U-Net/EfficientNet-B0, SegFormer-B0, DDRNet-23-Slim, and BiSeNetV2. With
PIDNet-S and ROD this yields nine model systems. Each is trained for the same
three seeds under both label policies, producing 54 scored runs.

FPN and U-Net are paired with both encoders to separate, within this limited
factorial subset, encoder and decoder effects. MobileNetV2 provides the lighter
convolutional encoder; EfficientNet-B0 provides a somewhat higher-capacity
compact encoder. SegFormer-B0 supplies a hierarchical transformer comparison.
DDRNet-23-Slim and BiSeNetV2 represent architectures designed specifically for
real-time segmentation.

Table~\ref{tab:model_implementation_details} records exact parameter counts
and initialisation. FPN, U-Net, and SegFormer are created through
\texttt{segmentation-models-pytorch}. DDRNet-23-Slim and BiSeNetV2 are adapted
from released PyTorch implementations under their licences. Every wrapper
returns one full-resolution two-class logits tensor so that the common loss and
metrics see the same interface.

\begin{table}[H]
\centering
\caption{Exact model systems, parameter counts, and initialisation in the
comparative study.}
\label{tab:model_implementation_details}
\small
\begin{tabularx}{\textwidth}{@{}lrlX@{}}
\toprule
\textbf{System} & \textbf{Parameters} & \textbf{Initialisation} & \textbf{Implementation note} \\
\midrule
PIDNet-S & 7,623,522 & Random & One final head; PAPPM and Light-Bag \\
ROD ViT-S & 29,106,050 & EfficientSAM ViT-S & 22.35M encoder parameters frozen \\
FPN/MobileNetV2 & 4,215,554 & ImageNet encoder & SMP FPN with \texttt{mobilenet\_v2} \\
FPN/EfficientNet-B0 & 5,759,614 & ImageNet encoder & SMP FPN with \texttt{efficientnet-b0} \\
U-Net/MobileNetV2 & 6,629,090 & ImageNet encoder & SMP U-Net with \texttt{mobilenet\_v2} \\
U-Net/EfficientNet-B0 & 6,251,614 & ImageNet encoder & SMP U-Net with \texttt{efficientnet-b0} \\
SegFormer-B0 & 3,714,658 & ImageNet MiT-B0 & SMP SegFormer decoder \\
DDRNet-23-Slim & 5,694,882 & Random & Two-resolution network; no stable pretrained checkpoint used \\
BiSeNetV2 & 3,341,202 & Released pretrained backbone & Auxiliary heads disabled for the common loss \\
\bottomrule
\end{tabularx}
\end{table}

Initialisation is not hidden because it is a major experimental factor on a
small dataset. PIDNet-S and DDRNet start from random weights; the other seven
systems use some form of pretraining, although ROD freezes its pretrained
encoder while the standard ImageNet encoders are trainable. The benchmark
therefore ranks the implemented recipes. It cannot isolate the causal effect
of architecture from pretraining without additional within-family
pretrained-versus-random experiments.

\section{Fixed-Budget Experimental Scope}

For every comparative run, the external input is $640\times384$, the
batch size is 4, the epoch budget is 15, and the model is selected by validation
mIoU. AdamW, the learning rate, cosine schedule, common CE+Dice objective,
policy-specific class weights, augmentation, 0.50 threshold, and three seeds
are fixed. Blue-only and blue+green experiments write to distinct processed
roots and output directories so that one policy cannot overwrite the other.

These controls support a restricted statement: under the tested CaT recipe,
model A has a higher three-seed mean than model B. They do not prove that the
ranking would remain under model-specific hyperparameter tuning, longer
training, more seeds, another dataset, or an embedded runtime. In particular,
the common 15-epoch budget may favour pretrained systems that converge quickly
and disadvantage randomly initialised models.

The nine-system experiment places the pairwise result in context. ROD's
advantage over PIDNet-S remains clear, but several pretrained
encoder--decoder systems achieve higher CaT overlap under the common recipe.
The results are therefore interpreted as a ranking of complete training
systems, with small differences read against their observed between-seed
variation.

\section{Convergence-Aware Blue+Green Follow-Up}

The 15-epoch matrix answers a matched-compute question, but the saved
validation histories show that it stops while several models are still
improving. A second experiment therefore repeats all nine architectures and
three seeds for the principal blue+green policy. It retains the same data,
external resolution, augmentation, objective, class weighting, optimiser,
initial learning rate, weight decay, threshold, and seed identities. Only the
training horizon and stopping rule change.

The follow-up sets 300 epochs as a ceiling, not as a presumed optimum. Cosine
decay runs from $2.5\times10^{-4}$ to $10^{-5}$ over the first 60 epochs; the
learning rate then remains at $10^{-5}$. Training cannot stop before epoch 60.
After that point it stops when validation mIoU has failed to improve by an
absolute increase of at least $10^{-4}$ on the 0-to-1 mIoU scale for 25
epochs. The checkpoint writer remains slightly more
sensitive than the stopping counter: it stores any strict validation-mIoU
increase, even if the increase is smaller than $10^{-4}$. Consequently, a
run may save its numerically highest checkpoint near the stopping epoch
without resetting the early-stopping counter.

This design gives every run time to pass the early learning phase while still
bounding unnecessary work after the validation curve has flattened. The
selected epoch is therefore called the \emph{validation-selected checkpoint
epoch}; it is not described as the universally optimal epoch for that
architecture. Model-specific objectives and schedules could still change both
the selected epoch and the final ranking.
